\label{ch:hawking_info}
% Source documents: 250, 325

\epigraph{The paradox does not resolve by adding new physics,
but by clarifying the premises.}{Doc.\,250}

\section{The Paradox and Its Three Premises}

Hawking showed in 1974 that black holes emit thermal radiation.
This radiation carries no information about the state of the infalling
matter --- it is purely thermal. If the black hole completely
evaporates, the original quantum information would be destroyed. That
contradicts the unitarity of quantum mechanics.

The paradox presupposes three premises, which do not all hold in FFGFT:

\begin{enumerate}
  \item Information is the complete description of the quantum state
    of all infalling particles --- an \emph{epistemic} category.
  \item Time is continuous and one-dimensionally directed.
  \item The classical singularity is physically real.
\end{enumerate}

FFGFT rejects all three --- not by assertion, but
through the geometric structure of the framework.

\section{Fundamental Information vs. Epistemic Information}

\begin{keyresult}
The fundamental information of a physical system in FFGFT is the
\textbf{topological charge} of its constituents: the winding numbers
$(n_\theta, n_\varphi, n_3, n_4)$ on $T^4$. These define \emph{which
kind} of particles exist --- not where they are, not in which exact
quantum state they are.
\end{keyresult}

Epistemic information --- what an observer can reconstruct ---
can be lost through thermalization. That is thermodynamic
irreversibility, not ontological information loss.

The Hawking paradox is an \emph{epistemic} problem: It asks whether an
observer outside the horizon can completely reconstruct the original quantum state.
FFGFT answers a different question:
Does the fundamental topological structure continue to exist?

\section{The $T^4$ Compactification Undermines Premise 2}

In FFGFT, time is a compact dimension of the $T^4$ torus, not a
continuous external parameter. The compactification means: There is
no absolute \enquote{forever inaccessible} --- the time coordinate is
periodic. What classically appears as irreversible information loss behind the
horizon is, in FFGFT, a geometric projection, not
ontological annihilation.

\section{The Singularity is Eliminated: Premise 3}

In FFGFT there is no singularity in the classical sense. The minimal
core radius $L_0 = \xi \cdot \lP$ prevents complete collapse. The
winding numbers remain conserved --- the topological charge is indestructible.
A black hole in FFGFT is not an information destroyer, but a
state of maximum winding density on the Schwarzschild scale.

\section{Hawking Temperature Without a Universal Bit Value --- Doc.\,325}

Doc.\,325 shows: The complete Hawking physics --- temperature,
spectrum, information conservation --- follows from $\tilde{T} \cdot m = 1$,
the $\mathbb{Z}_3$ structure and the system-dependent bit energy
$\Ebit = \hbar c / L$, \emph{without} a universal bit value.

That is the decisive difference from Matzke: Matzke sets
$m_{\text{bit}} = m_P \ln 2 / (2\pi)$ as a universal constant and
derives the Schwarzschild radius from it. FFGFT arrives at the same
numerical results --- but without this stipulation. The stability
threshold $n_{\text{thresh}} = 6{,}41$ in Matzke depends on the bit value parametrization,
not on the temperature physics.

\begin{keyresult}
Matzke and FFGFT converge on \textbf{fermion generation = fundamental
algebraic unit}. They diverge on \textbf{bit values}: universal
(Matzke, Landauer at $T_P$) vs.\ system-dependent (FFGFT, $\Ebit = \hbar c/L$).
This difference concerns the stability threshold and the mass of a bit.
The Hawking temperature formula itself is unaffected: Both routes give
the same number, but only the FFGFT route manages without a universal bit value.
\end{keyresult}